\section{Statics}
\label{sec:expr-proc}

% Figures in this section: syntax of expressions (\cref{fig:syntax-expressions}),
% run-time syntax (\cref{fig:syntax-runtime}), reduction (expressions:
% \cref{fig:expression-reduction}, processes: \cref{fig:process-reduction}),
% process congruence (\cref{fig:configuration-congruence}).

\begin{figure}
  \begin{align*}
    \shortintertext{\textbf{Types}}
    \Mob \grmeq
              & \MobStationary \grmor \MobMobile
    \tag{Mobilities}                                                                           \\
    \Dir \grmeq
              & \DirLeft \grmor \DirRight \grmor \DirUnord \grmor \DirNone
    % \grmor \DirNone
    \tag{Directions}                                                                           \\
    \Eff \grmeq
              & \EffPure \grmor \EffImpure
    \tag{Effects with $\EffPure < \EffImpure$}                                                 \\
    \T \grmeq & \BUnit
    \grmor \TArr{\Mob}{\Dir,\Eff}\T\T
    \grmor \T \TPair \T
    \grmor \TVariant{\ell : \T[\ell]}
    \grmor \S
    \tag{Types}                                                                                \\
    \S \grmeq
              & \SSkip \grmor \SSeq\S\S \grmor \STerm \grmor \SWait \grmor
    \SSend\T \grmor \SRecv\T
    \\
              &
    \grmor \SChoice{\ell:\S[\ell]}{\ell\in L}
    \grmor \SBranch{\ell:\S[\ell]}{\ell\in L}
    \\
              &
    \grmor \SRet \grmor \SAcq
    \grmor
    \SRec\SVar\S \grmor \SVar % no rec yet
    \tag{Session types}
    \\
    \Ctx \grmeq
              & \CtxEmpty \grmor \CtxVar \EVar\T \grmor \Ctx\CtxSeq\Ctx \grmor \Ctx\CtxPar\Ctx
    \tag{Typing environments}                                                                  \\
    \shortintertext{\textbf{Expressions}}
    \C \grmeq
              & \CUnit
    \grmor \CNew
    \grmor \CFork
    \grmor \CTerm
    \grmor \CWait
    \grmor \CSend
    \grmor \CRecv
    \\\grmor &
      \CSplit \grmor \CSSplit
      \grmor \CDrop \grmor \CAcq
      \grmor \CSelect k
      \grmor \CBranch
    \tag{Constants}                                                                            \\
    \E \grmeq
              & \C \grmor \EVar \grmor \EAbs\EVar\E \grmor \EAppD{\E}{\E}{\Dir}
    \grmor \ESeq\E\E
    \grmor \EPair\E\E
    \grmor \ELetPair{\EVar}{\EVar}{\E}{\E}                                                     \\
    \grmor    & \EInj i \E
    %\grmor \ECase\E{\EVar}{\E}{\EVar}{\E}
    \grmor \ECase \E {\ECaseBranch* i \EVar \E}
    \grmor \EMu[\EVar] \E
    % & \grmor \EMatch{\E}{\ell:\E_\ell}{\ell\in L}
    % \grmor \EFork\E \grmor \ELet\EVar\E\E
    \tag{Expressions}
    \\
    \shortintertext{\textbf{Processes}}
    \BindGroup \grmeq
              & (\EVar \grmor \BSep)^*
    \tag{Binder groups}
    \\
    \Bind \grmeq
              & \EVar^*
    \tag{Variable sequences}
    \\
    \Proc \grmeq
              & \PExp\E
    \grmor \PPar\Proc\Proc
    \grmor \PNew{\BindGroup}{\BindGroup}{\Proc}
    \tag{Processes}
  \end{align*}
  \caption{Syntax of types, expressions, processes}
  \label{fig:syntax}
\end{figure}

Figure \cref{fig:syntax} collects the source-language syntax used throughout
this paper. The figure is organized into three groups: types, expressions, and
processes.

The type group introduces four orthogonal classifiers. Mobilities
($\MobStationary$, $\MobMobile$) record whether values are stationary or 
can be moved across process boundaries. Directions ($\DirLeft$,
$\DirRight$, $\DirUnord$, $\DirNone$) govern evaluation order and context
composition in function and pair constructs. An effect is either pure or impure
and together they form a two-point lattice with $\EffPure < \EffImpure$. Types
$\T$ combine ordinary functional forms (unit, arrows, pairs, variants) with
session types $\S$. Session types use sequential composition as their central
connective and include communication, branching/selection, and recursion.
Typing environments $\Ctx$ combine variable bindings using ordered and
unordered composition.

The expression group separates constants $\C$ from compound expressions $\E$.
Constants cover channel creation, process creation, communication primitives,
split operations, and synchronization primitives. They are better understood
together with their type in the next section.

Expressions then add variables, abstraction and application, sequencing,
products, sums, and recursive terms. Following our directional presentation,
application carries an explicit direction annotation. We distinguish left
($\DirLeft$) and right ($\DirRight$) application syntactically, but unlike
previous work~\cite{saffrich2024lawordertypestateborrowing, DBLP:journals/pacmpl/SaffrichS0V25} we do not distinguish left and right abstraction. The
direction determines both the evaluation order and context position in the body
of the function. Throughout this paper, omitted parentheses and omitted
direction annotations are understood as standard, right-directed
CBV-application.

Binder groups $\BindGroup$ are flat lists of channel names $\EVar$ and
separators $\BSep$. The metavariable $\Bind$ ranges over separator-free
variable sequences. We write $\BEmp$ for the empty list and
$\BWithSep{\Bind}{\BindGroup}$ for the concatenation of $\Bind$, $\BSep$, and
$\BindGroup$. The paper only creates groups where two separators never occur
side-by-side without an intervening channel name, but the syntax does not
enforce this invariant. Processes $\Proc$ include a single expression, parallel
composition of processes, and channel introduction with two binder groups.

Binder sequence concatenation is written by juxtaposition with unit $\BEmp$.

%%% Local Variables:
%%% mode: latex
%%% TeX-master: "../popl2027.tex"
%%% End:
